\documentclass[11pt]{article}
\usepackage[T1]{fontenc}
\usepackage[utf8]{inputenc}
\usepackage{lmodern}
\usepackage{amsmath,amssymb}
\usepackage{longtable,booktabs,array}
% Compatibility with Pandoc's captionless tables on older longtable versions.
\ifcsname c@none\endcsname\else\newcounter{none}\fi
\usepackage{graphicx}
\usepackage[margin=25mm]{geometry}
\usepackage{hyperref}
\hypersetup{colorlinks=true,linkcolor=blue,urlcolor=blue}
\providecommand{\tightlist}{\setlength{\itemsep}{0pt}\setlength{\parskip}{0pt}}
\providecommand{\pandocbounded}[1]{#1}
\setlength{\emergencystretch}{3em}
\setcounter{secnumdepth}{-1}
\begin{document}
\section{Serre's comparison theorems and Chow's
theorem}\label{serres-comparison-theorems-and-chows-theorem}

\emph{Written and self-checked by GPT-6.1 Sol (OpenAI) in Codex, Ultra
setting, October 2026. Analytic prerequisite integration by GPT-6 Astra
(OpenAI), Codex, Ultra. The linked analytic bridge was written and
self-checked by Claude Opus 5.5; a full independent review is not
claimed. Independently authored material: CC0.}

Locally, analytification preserves exact sequences but adds holomorphic
functions. On a projective complex scheme, the three comparison theorems
say much more: coherent cohomology agrees, every analytic morphism of
coherent sheaves is algebraic, and every coherent analytic sheaf is
algebraic. We prove these assertions in that order, including schemes
with nilpotents. The proof of the last assertion needs a
finite-dimensional analytic cohomology theorem; using the comparison
theorem to assert finiteness for an arbitrary analytic sheaf at this
stage would be circular.

Throughout, schemes are separated and of finite type over \(\mathbf C\).
A variety means a reduced such scheme; a complex space may have
nilpotents. Analytification retains the full defining ideal, as
constructed in
\href{complex-analytic-spaces-and-analytification.html}{Complex analytic
spaces and analytification}. Write \(P=\mathbf P^n_{\mathbf C}\),
\(P^{\mathrm{an}}=\mathbf P^n(\mathbf C)\), and
\(\mathcal O(d)^{\mathrm{an}}\) for the analytification of a twist. We
use the generation, vanishing and projective cohomology proved in
\href{serres-theorems-on-projective-schemes.html}{Serre's theorems} and
\href{cohomology-of-projective-space.html}{Cohomology of projective
space}.

\subsection{1. The analytic cohomology
foundation}\label{the-analytic-cohomology-foundation}

Two analytic theorems are required for the cohomology arguments below.
Their programme proofs are in \emph{Complex analytic spaces and coherent
sheaves}. Read the vanishing and finiteness lessons before this lesson:
neither obtains analytic cohomology from GAGA or from algebraization.
Jean-Pierre Demailly's freely available
\href{https://www-fourier.univ-grenoble-alpes.fr/~demailly/manuscripts/agbook.pdf}{\emph{Complex
Analytic and Differential Geometry}} is further reading. The
\href{prerequisites.html}{prerequisite and proof guide} identifies the
exact statements, proof locations and remaining dependencies.

\textbf{Analytic vanishing theorem.} If a complex manifold admits a
smooth strictly plurisubharmonic exhaustion, every coherent analytic
sheaf on it has zero cohomology in positive degrees. An exhaustion has
relatively compact sublevel sets; strict plurisubharmonicity means its
complex Hessian is positive definite.
\href{../complex-analytic-spaces-and-coherent-sheaves/theorems-a-and-b-on-stein-manifolds.html\#4-theorem-b}{Theorems
A and B on Stein manifolds, Theorem 4.1}, supplies the proof through
local coherent resolutions, approximation and exhaustion. The assertion
applies to every coherent sheaf, not only vector bundles; it gives
vanishing in every positive degree, not merely finite dimension.

\textbf{Analytic finiteness theorem (Cartan--Serre).} On a compact
complex analytic space, the cohomology of every coherent analytic sheaf
is finite-dimensional in every degree.
\href{../complex-analytic-spaces-and-coherent-sheaves/finiteness-on-compact-complex-spaces.html\#2-the-finiteness-theorem}{Finiteness
on compact complex spaces, Theorem 2.1}, gives the proof using finite
nested acyclic covers. Restriction of coherent sections from a larger
open set to a relatively compact smaller one is a compact operator, by
the coherent version of Montel's theorem. The induced map between their
Čech complexes is an isomorphism on cohomology. The compact-operator
criterion is developed in the preceding lesson,
\href{../complex-analytic-spaces-and-coherent-sheaves/frechet-spaces-of-sections-and-schwartzs-theorem.html}{Fréchet
spaces of sections and Schwartz's theorem}. It yields finite dimension
and closed coboundaries. These arguments apply to coherent analytic
sheaves before any algebraization is known.

For our standard projective cover \(U_i=\{T_i\ne0\}\), every nonempty
intersection is biholomorphic to

\[
(\mathbf C^*)^r\times\mathbf C^{n-r}.
\]

The function

\[
\psi(z)=\sum_{j=1}^n|z_j|^2+\sum_{j=1}^r|z_j|^{-2}
\]

is a strictly plurisubharmonic exhaustion there. The first sum makes the
Hessian positive, while the second prevents escape toward a deleted
coordinate hyperplane. The vanishing theorem therefore makes this cover
acyclic for every coherent analytic sheaf. By the Čech comparison proved
earlier in the course, its ordered complex computes cohomology. There
are \(n+1\) opens, so

\[
H^q(P^{\mathrm{an}},M)=0\qquad(q>n)
\tag{1}
\]

for every coherent analytic module \(M\). This bound concerns all
coherent analytic modules, including those not yet known to be
algebraic.

\subsection{2. Starting the comparison: the structure sheaf and
twists}\label{starting-the-comparison-the-structure-sheaf-and-twists}

\textbf{Lemma 2.1.} The analytic structure sheaf of projective space has
cohomology \(\mathbf C\) in degree zero and zero in every positive
degree. The natural comparison from algebraic cohomology is an
isomorphism.

\textbf{Proof.} For \(n=0\), the space is a point. For general \(n\),
use the acyclic cover of Section 1. Put
\(\Omega_S=\{T\in\mathbf C^{n+1}:T_i\ne0\text{ for }i\in S\}\). A
function on \(U_S=\bigcap_{i\in S}U_i\) pulls back to a homogeneous
holomorphic function of weight zero on this product of punctured planes
and planes. Conversely a weight-zero function descends by setting any
\(T_j=1\), \(j\in S\); homogeneity makes this independent of the
representative. Work directly on \(\Omega_S\), without choosing
projective ratios for the expansion.
\href{../../human/elementary-analysis/laurent-series-and-homogeneous-projections.html\#2-products-of-annuli-and-discs}{Laurent
series and homogeneous projections, Theorems 1--4}, proves the unique,
normally convergent product expansion and its homogeneous decomposition.
The allowed exponents are

\[
\alpha\in\mathbf Z^{n+1},\qquad
\sum_i\alpha_i=0,\qquad
\alpha_j\geq0\quad(j\notin S).
\]

Indeed uniqueness applied to \(f(2T)=f(T)\) makes the coefficient of
\(T^\alpha\) vanish unless its total exponent is zero. Group the
expansion by \(N=\{j:\alpha_j<0\}\). There are only finitely many such
sets.
\href{../../human/elementary-analysis/laurent-series-and-homogeneous-projections.html\#3-extending-the-parts-with-specified-negative-powers}{Theorem
3 of the Laurent unit} proves that each part extends to the whole
\(\Omega_N\): on any compact subset, choose smaller positive contour
radii in the negative-power coordinates and larger contour radii in the
others. The product geometric estimate proves normal convergence there,
including along the newly allowed zero coordinates. The extension still
has weight zero. Uniqueness of coefficients on \((\mathbf C^*)^{n+1}\)
makes every such projection compatible with restriction between chart
intersections. Thus each fixed \(N\) has one coefficient space,
independent of the containing face \(S\), rather than a separate choice
for every chart.

For fixed \(N\), the allowed faces of the Čech complex are exactly the
nonempty \(S\) containing \(N\). If \(N\) is proper and nonempty, choose
\(j\notin N\). Insert \(j\) with the sign of its ordered position,
extending the component on \(\Omega_{S\cup\{j\}}\) first to
\(\Omega_N\), then restricting to \(\Omega_S\). If \(j\) is already in
the target face, set the contraction to zero.
\href{../../human/elementary-analysis/laurent-series-and-homogeneous-projections.html\#4-homogeneous-functions-and-the-ordered-complex}{Theorem
5} gives the signed formula and proves \(dh+hd=1\), including degree
zero. The term removing the inserted \(j\) is the identity; every other
pair differs by one in its insertion/removal sign and cancels. This is
an operation on normally convergent holomorphic components, not only a
formal monomial calculation.

If \(N\) is empty, all exponents are nonnegative and sum to zero, so the
only monomial is 1. This is the ordinary simplex complex of constants,
with cohomology \(\mathbf C\) in degree zero. If \(N\) were all indices,
their sum would be negative, so this part is absent. These contractions
prove the analytic assertion. Algebraically the same cohomology was
computed in our projective-space lesson, and comparison sends a constant
to the identical constant. \(\square\)

For any algebraic module \(F\), the map on global sections to those of
\(F^{\mathrm{an}}\) extends canonically to maps

\[
c_F^q:H^q(X,F)\longrightarrow H^q(X^{\mathrm{an}},F^{\mathrm{an}}).
\tag{2}
\]

Indeed analytification is an exact functor on all modules, by local
flatness. The target groups therefore form a cohomological delta
functor, and the universal derived functors of algebraic global sections
extend the degree-zero map uniquely. Consequently (2) is natural and
commutes with connecting maps in long exact sequences. For coherent
modules, it also commutes with closed-immersion pushforward, whose
stalkwise quotient description was proved in the preceding lesson.

\textbf{Lemma 2.2.} Comparison is an isomorphism for \(\mathcal O(d)\)
on \(\mathbf P^n\), for every integer \(d\) and every degree.

\textbf{Proof.} Induct on \(n\). On \(\mathbf P^0\) every twist is a
one-dimensional module, including negative twists. For \(n>0\), let
\(i:E\hookrightarrow P\) be a hyperplane. Its equation gives

\[
0\to\mathcal O(d-1)\to\mathcal O(d)
\to i_*\mathcal O_E(d)\to0.
\tag{3}
\]

Analytification preserves this sequence. The induction hypothesis
supplies comparison on \(E\). The two long exact sequences show that
comparison for \(\mathcal O(d)\) in all degrees is equivalent to
comparison for \(\mathcal O(d-1)\): for either unknown term, take five
successive terms with that term in the middle; all four other comparison
maps are isomorphisms. Begin with \(d=0\), proved in Lemma 2.1, and move
upward and downward. \(\square\)

\subsection{3. The first comparison
theorem}\label{the-first-comparison-theorem}

\textbf{Theorem 3.1 (cohomological GAGA).} If \(X\) is projective over
\(\mathbf C\) and \(F\) is coherent, (2) is an isomorphism for every
\(q\geq0\).

\textbf{Proof.} Embed \(X\) as a closed subscheme of \(P\), with its
full coherent defining ideal. Closed-immersion pushforward is exact,
preserves coherence and cohomology, and commutes with analytification:
locally a module over a quotient ring is viewed as a module over the
ambient ring, and both analytic constructions tensor its presentation
with the same flat analytic ring. None of these operations replaces the
ideal by its radical. Thus it suffices to prove the result on \(P\).

Descend on \(q\), simultaneously for all coherent sheaves. For \(q>n\),
both groups vanish: algebraically by the standard affine cover,
analytically by (1). Choose a presentation

\[
0\to R\to L\to F\to0,
\tag{4}
\]

where \(L\) is a finite sum of twists and \(R\) is coherent. Suppose
comparison in degree \(q+1\) is already known for every coherent sheaf.
Lemma 2.2 gives comparison for \(L\) in every degree.

To prove surjectivity for \(F\) in degree \(q\), start with an analytic
class. Its boundary in \(H^{q+1}(R^{\mathrm{an}})\) lifts uniquely to an
algebraic class. That lift maps to zero in \(H^{q+1}(L)\), since its
analytic image does and comparison for \(L\) is injective. Exactness
lifts it to \(H^q(F)\). The difference between this class's analytic
image and the original class has zero boundary, so comes from
\(H^q(L^{\mathrm{an}})\). Comparison for \(L\) lifts the difference,
establishing surjectivity.

This proves surjectivity in degree \(q\) for every coherent sheaf, in
particular for \(R\). If an algebraic class in \(H^q(F)\) has zero
analytic image, its boundary is zero by injectivity in degree \(q+1\)
for \(R\), so it comes from \(H^q(L)\). That lift maps analytically into
the image of \(H^q(R^{\mathrm{an}})\). Surjectivity for \(R\) lifts this
correcting class. Subtract it in \(H^q(L)\); comparison for \(L\) forces
the corrected lift to be zero. Thus the original class was zero. The
induction proves injectivity and surjectivity in all degrees.
\(\square\)

The simultaneous induction is essential: we first establish surjectivity
for all coherent sheaves before using it for the kernel \(R\). No finite
global locally free resolution of \(F\) has been assumed.

\subsection{4. The second comparison
theorem}\label{the-second-comparison-theorem}

\textbf{Theorem 4.1 (full faithfulness).} For coherent \(F,G\) on a
projective complex scheme,

\[
\operatorname{Hom}_X(F,G)
\xrightarrow{\sim}
\operatorname{Hom}_{X^{\mathrm{an}}}(F^{\mathrm{an}},G^{\mathrm{an}}).
\]

\textbf{Proof.} The internal Hom \(A=\mathcal Hom(F,G)\) is coherent.
Finite presentations and local flatness give

\[
A^{\mathrm{an}}\cong\mathcal Hom(F^{\mathrm{an}},G^{\mathrm{an}}),
\]

as proved in Proposition 5.1 of the preceding lesson. Apply Theorem 3.1
in degree zero to \(A\). These global sections are exactly the two
morphism groups, and the comparison takes a morphism to its
analytification. \(\square\)

In particular, an analytic isomorphism between algebraic coherent
sheaves lifts uniquely to an algebraic isomorphism: lift it and its
inverse, then use faithfulness to identify their compositions with the
identities.

\subsection{5. The third comparison
theorem}\label{the-third-comparison-theorem}

\textbf{Lemma 5.1 (analytic generation).} Every coherent analytic module
\(M\) on \(P^{\mathrm{an}}\) has \(M(d)\) generated by finitely many
global sections for all sufficiently large \(d\).

\textbf{Proof.} We prove this together with essential surjectivity,
inducting on \(n\). At \(n=0\) a coherent module is a finite-dimensional
vector space. Assume essential surjectivity on every hyperplane
\(E\cong\mathbf P^{n-1}\). On \(E\), Theorem 3.1 and algebraic Serre
vanishing now imply vanishing of all positive cohomology of sufficiently
large twists of every coherent analytic module.

Fix \(x\in P^{\mathrm{an}}\) and choose a hyperplane through \(x\), with
equation \(t\). Multiplication has coherent kernel and cokernel:

\[
0\to C\to M(-1)\xrightarrow{t}M\to B\to0.
\]

Both \(C\) and \(B\) are annihilated by \(t\), so are coherent modules
on \(E^{\mathrm{an}}\). This remains true when multiplication by \(t\)
is not injective. Let \(L_d\) be the image of \(M(d-1)\to M(d)\). The
two short exact sequences give, for large \(d\), surjections

\[
H^1(M(d-1))\twoheadrightarrow H^1(L_d)
\twoheadrightarrow H^1(M(d)),
\tag{5}
\]

because \(H^2(C(d))=H^1(B(d))=0\). By the Cartan--Serre finiteness
theorem in Section 1, these dimensions are finite. Thus the nonnegative
integers \(h^1(M(d))\) eventually stop decreasing. In the stable range
both arrows in (5) are isomorphisms. The second short exact sequence
then shows that

\[
H^0(M(d))\twoheadrightarrow H^0(B(d)).
\tag{6}
\]

By induction \(B\) is algebraic on \(E\); its large twists are generated
by global sections. Lift generators using (6). At \(x\), these generate
\(M(d)_x/tM(d)_x\); Nakayama makes them generate \(M(d)_x\), since \(t\)
belongs to the local maximal ideal.

Generation by finitely many fixed sections holds on a neighbourhood,
because their coherent cokernel vanishes there. Once generation holds at
a point in degree \(d\), multiplying these sections by a homogeneous
coordinate nonzero at that point proves generation in every degree
\(d'\geq d\). Compactness supplies finitely many neighbourhoods and a
common bound on \(d\). At each such degree take the union of their
finite generating families. This proves the lemma. \(\square\)

\textbf{Theorem 5.2 (essential surjectivity).} Every coherent analytic
module on \(X^{\mathrm{an}}\), for projective \(X\), is isomorphic to
\(F^{\mathrm{an}}\) for some coherent algebraic \(F\). Together with
Theorem 4.1, analytification is an equivalence of coherent categories.

\textbf{Proof.} First take \(X=P\), continuing the induction of Lemma
5.1. Global generation gives a surjection \(L_0^{\mathrm{an}}\to M\),
where \(L_0\) is a finite sum of equal negative twists. Its kernel \(R\)
is coherent. Generate a twist of \(R\) as well, to obtain

\[
L_1^{\mathrm{an}}\xrightarrow{v}L_0^{\mathrm{an}}\to M\to0.
\]

Full faithfulness lifts \(v\) uniquely to an algebraic map
\(u:L_1\to L_0\). Put \(F=\operatorname{coker}u\). Exactness of
analytification gives \(F^{\mathrm{an}}\cong M\), completing the
induction on \(n\).

For \(i:X\hookrightarrow P\), algebraize \(i_*M\) to \(G\) on \(P\). If
\(I\) is the ideal of \(X\), then
\((IG)^{\mathrm{an}}=I^{\mathrm{an}}G^{\mathrm{an}}=0\). Exactness and
faithfulness give \(IG=0\). Hence \(G=i_*F\) for a unique coherent
module \(F\) on \(X\). Compatibility with closed-immersion pushforward
identifies \(F^{\mathrm{an}}\) with \(M\). \(\square\)

The correct uniqueness assertion fixes the identification. If
\((F,\phi)\) and \((G,\psi)\) are algebraizations equipped with
isomorphisms to \(M\), there is a unique isomorphism \(u:F\to G\)
satisfying \(\psi\circ u^{\mathrm{an}}=\phi\). Without specified
identifications, an algebraic sheaf can have many automorphisms: every
nonzero scalar acts on \(\mathcal O_X\). Thus its bare isomorphism class
is unique, but its possible isomorphisms are not.

\subsection{6. Chow's theorem and geometric
consequences}\label{chows-theorem-and-geometric-consequences}

\textbf{Theorem 6.1 (Chow).} Every closed analytic subset of a
projective complex scheme is the analytification of a unique reduced
closed algebraic subscheme.

\textbf{Proof.} Its analytic vanishing ideal
\(J\subset\mathcal O_{X^{\mathrm{an}}}\) is coherent by Cartan
coherence. Algebraize \(J\) as a coherent module by Theorem 5.2, and
lift its inclusion into the structure sheaf by Theorem 4.1. Exactness
and faithfulness make the lift injective, with image an algebraic
coherent ideal \(I\) whose analytification is \(J\).

The ideal \(I\) is radical. At every closed stalk, if \(a^m\in I_x\),
the image of \(a\) belongs to the radical ideal
\(J_x=I_x\mathcal O_{X^{\mathrm{an}},x}\); faithful flatness contracts
this to \(a\in I_x\). Radicality at closed stalks suffices: a nonzero
nilpotent ideal in the coherent quotient would have support containing a
closed point. The reduced closed scheme defined by \(I\) therefore
analytifies to the given analytic subset. If two algebraic ideals have
the same analytic extension, faithful flatness gives equal closed
stalks, and coherent support detects equality globally. This proves
uniqueness. \(\square\)

\textbf{Theorem 6.2.} Every holomorphic morphism between the
analytifications of projective complex schemes is the analytification of
a unique algebraic morphism.

\textbf{Proof.} Analytification commutes with products, by the
polynomial chart and quotient constructions in the preceding lesson. The
graph is a closed complex subspace of \((X\times Y)^{\mathrm{an}}\):
locally its ideal is generated by the target coordinate functions minus
their holomorphic pullbacks, together with the target equations. Its
structure sheaf is that of \(X^{\mathrm{an}}\), including any
nilpotents. Its ideal \(J\) is coherent. Theorem 5.2 algebraizes \(J\),
and Theorem 4.1 algebraizes its inclusion in the structure sheaf.
Exactness and faithfulness make the image a coherent algebraic ideal
\(I\) with \(I^{\mathrm{an}}=J\). This gives a closed subscheme
\(Z\subset X\times Y\) whose full analytification is the graph; reduced
Chow alone would lose its infinitesimal structure.

The first projection \(p:Z\to X\) is proper and its analytification is
an isomorphism. Each closed fiber has one complex point and is
zero-dimensional. Upper semicontinuity of fiber dimension for proper
maps shows that all fibers are zero-dimensional: any nonempty closed
locus of positive-dimensional fibers would contain a closed point. Thus
\(p\) has finite fibers and is finite by
\href{the-theorem-on-formal-functions.html\#5-dimension-and-finiteness-consequences}{Formal
functions, Theorem 5.2}.

At a closed point \(x\), the finite algebra \((p_*\mathcal O_Z)_x=B\)
over \(A=\mathcal O_{X,x}\) has only one maximal ideal, since the closed
fiber has one point. Its maximal-ideal topology is cofinal with its
\(\mathfrak m_A\)-adic topology. Thus
\(B\otimes_A\widehat A=\widehat B\). Local analytic comparison and the
analytic graph isomorphism identify the map \(\widehat A\to\widehat B\)
with an isomorphism. Faithful flatness of completion makes \(A\to B\) an
isomorphism. The kernel and cokernel of
\(\mathcal O_X\to p_*\mathcal O_Z\) are coherent and vanish at every
closed point, so they vanish everywhere. A finite map is the relative
spectrum of this algebra, hence \(p\) is an isomorphism. The second
projection composed with \(p^{-1}\) is the required morphism. Two
algebraic graph ideals with the same analytic extension agree by
faithful flatness at closed points and coherent support. Their first
projections and second projections therefore give the same algebraic
morphism, proving uniqueness. \(\square\)

\textbf{Proposition 6.3.} For \(n\geq1\), every analytic line bundle on
\(\mathbf P^n(\mathbf C)\) is \(\mathcal O(d)^{\mathrm{an}}\) for a
unique \(d\in\mathbf Z\).

\textbf{Proof.} GAGA gives a coherent algebraic module \(L\). It is a
line bundle: at a closed point choose one element lifting an analytic
fiber basis. Nakayama and faithful flatness make the resulting map
\(A\to L_x\) onto; after analytic tensor it is a map between rank-one
free modules with nonzero residue, hence an isomorphism. Faithfulness
makes its algebraic kernel zero. The locally free locus is open, and its
closed complement contains no closed point, so is empty.

Here is the divisor argument, including why it classifies the bundle.
Let \(K\) be the rational function field of \(P\) and choose a nonzero
rational section \(s\) of \(L\). In a local frame \(e_a\), write
\(s=g_a e_a\), with \(g_a\in K^*\). On overlaps the quotients
\(g_a/g_b\) are units. Thus the orders of the \(g_a\) along irreducible
hypersurfaces define a finite divisor \(D\); finiteness follows from a
finite trivializing cover and from the finite numerator and denominator
factorizations of each \(g_a\). The sheaf \(\mathcal O(D)\) consists
locally of rational functions \(h\) such that \(hg_a\) is regular, and
\(h\mapsto hs\) identifies it with \(L\).

The affine charts of \(P\) are polynomial rings, and their local rings
are factorial. One can see the polynomial assertion inductively from
Gauss's argument: in a factorial ring the product of two primitive
polynomials is primitive, since reduction modulo each prime factor gives
a product of two nonzero polynomials over a domain; extract the greatest
common divisor of the coefficients and factor over the fraction field.
Clearing denominators and this primitive-product assertion give
existence and uniqueness of polynomial factorization. Localization
simply discards factors made invertible. Consequently the orders along
prime hypersurfaces determine a local Cartier divisor: a rational
function with zero order at every prime of a factorial local ring is a
unit, by cancelling its numerator and denominator factors.

Every prime hypersurface of \(P\) has an irreducible homogeneous
equation \(f\). To justify the assertion, its cone has a homogeneous
prime ideal of height one: on any nonempty coordinate chart the
punctured cone is the hypersurface times \(\mathbf G_m\), so its
codimension is one. In a factorial ring a height-one prime is principal:
factor a nonzero element of the prime, choose an irreducible factor in
it, and use the chain \((0)\subset(f)\subseteq\mathfrak p\) and height
one. A generator of a homogeneous principal ideal is homogeneous. Indeed
take its highest and lowest nonzero homogeneous components; divisibility
of either component by the generator forces the two degrees to agree,
because highest and lowest degrees add under multiplication in a
polynomial domain.

If \(\deg f=e\) and \(H=V(T_0)\), the degree-zero rational function
\(f/T_0^e\) has divisor \(V(f)-eH\): in a local frame the numerator has
order one on its prime hypersurface and the denominator has order \(e\)
on \(H\). For \(D=\sum m_iV(f_i)\), multiply these rational functions to
see that \(D\) differs from \(dH\), \(d=\sum m_i\deg f_i\), by a
principal divisor. Multiplication by that rational function gives an
isomorphism \(\mathcal O(D)\cong\mathcal O(dH)=\mathcal O(d)\). The last
equality follows directly from the transition functions \((T_b/T_a)^d\)
on the standard charts. Therefore \(L\cong\mathcal O(d)\).

If two integers give isomorphic twists, their difference gives a trivial
twist. Its global-section dimension would be one. Our projective-space
calculation gives zero for a negative degree and \(\binom{d+n}{n}\) for
a nonnegative degree, which is one only for \(d=0\) when \(n\geq1\).
This proves uniqueness. On \(\mathbf P^0\), there is just the trivial
line bundle and every twist represents it. \(\square\)

A compact Riemann surface holomorphically embedded in projective space
consequently has algebraic image. It is an algebraic smooth projective
curve: dimensions and regularity agree at complex points by local
comparison. The hypothesis includes an embedding; this conclusion does
not assert that an arbitrary compact analytic space is projectively
embeddable.

\subsubsection{From a projective modification to a proper
scheme}\label{from-a-projective-modification-to-a-proper-scheme}

\textbf{Lemma 6.4.} A reduced proper complex scheme \(X\) admits a
projective surjection \(p:Y\to X\), with \(Y\) reduced and projective,
which is an isomorphism over an open \(U\) containing every generic
point of \(X\). Every proper complex scheme has compact analytification.

\textbf{Proof.} For each reduced irreducible component \(X_i\), use the
Chow construction proved in
\href{proper-morphisms-and-coherent-direct-images.html\#4-a-projective-modification-supplies-the-witness}{Coherence
of higher direct images under proper morphisms, Lemma 4.0}. It gives a
projective modification \(Y_i\to X_i\), an isomorphism over a nonempty
open, with \(Y_i\) integral and projective over \(\mathbf C\). Set
\(Y=\coprod_iY_i\). A finite disjoint union is projective: embed the
ambient projective spaces as pairwise disjoint coordinate subspaces of
one larger projective space. The same construction over \(X\) makes
\(p\) projective. Remove the exceptional closed subsets and the
intersections between distinct components. The remaining open \(U\)
contains every generic point, and \(p^{-1}U\to U\) is an isomorphism.

Every nonempty finite-type fibre over a complex point has a complex
point, by the earlier algebraic Nullstellensatz, so \(p^{\mathrm{an}}\)
is surjective. The space \(Y^{\mathrm{an}}\) is compact: a projective
closed embedding identifies it with a closed subspace of complex
projective space, the continuous image of the unit sphere. Hence
\(X^{\mathrm{an}}\) is compact. An arbitrary \(X\) and its reduction
have the same analytic underlying space by the full quotient
construction, so compactness follows in that case too. \(\square\)

\subsubsection{Comparing projective direct
images}\label{comparing-projective-direct-images}

This comparison is relative; its base need not be proper. It precedes
the use of a projective modification over a nonprojective proper scheme.

\textbf{Lemma 6.5 (twists with holomorphic parameters).} Let
\(D\subset\mathbf C^m\) be a polydisc and let
\(r:\mathbf P^n(\mathbf C)\times D\to D\) be the projection. For every
\(d\in\mathbf Z\), \[
H^q(\mathbf P^n(\mathbf C)\times D,\mathcal O(d))
=V_d^q\otimes_{\mathbf C}\mathcal O(D),
\tag{7}
\] where \(V_d^q=H^q(\mathbf P^n_{\mathbf C},\mathcal O(d))\), with the
same monomial bases as in the algebraic projective-space computation.
These identifications commute with restriction in \(D\) and with the
comparison maps.

\textbf{Proof.} Use the \(n+1\) standard projective charts, multiplied
by \(D\). Every finite intersection is
\((\mathbf C^*)^a\times\mathbf C^{n-a}\times D\). It has a smooth
strictly plurisubharmonic exhaustion: add \(\sum|w_j|^2\),
\(\sum_{j\le a}|w_j|^{-2}\), and \(\sum_k(R_k^2-|z_k-b_k|^2)^{-1}\),
where \(D=\prod_k\{|z_k-b_k|<R_k\}\). The first and last terms give a
positive complex Hessian, and these terms tend to infinity on every
escape from a compact subset. Cartan's Theorem B on Stein manifolds
therefore makes this finite cover acyclic for every coherent module, in
particular each twist. Its ordered Čech complex computes the cohomology.

For a nonempty set \(S\subset\{0,\ldots,n\}\), put \[
\Omega_S=\{T\in\mathbf C^{n+1}:T_i\ne0\ (i\in S)\}.
\] A section of \(\mathcal O(d)\) on the corresponding intersection,
pulled back to \(\Omega_S\times D\), is a holomorphic function of
homogeneous weight \(d\) in \(T\). Its product Laurent expansion is \[
\sum_{\substack{\alpha\in\mathbf Z^{n+1}\\
\sum\alpha_i=d,\ \alpha_j\ge0\ (j\notin S)}}
c_\alpha(z)T^\alpha .
\tag{8}
\] Indeed multitorus Cauchy integrals give the coefficients; the
integrands are jointly holomorphic, so the coefficients are holomorphic
in \(z\). Uniqueness and the identity \(f(2T,z)=2^df(T,z)\) impose the
weight condition. The same Cauchy estimates are uniform for \(z\) in a
compact subset of \(D\); hence this Laurent series converges normally on
every compact subset of \(\Omega_S\times D\).

Group its terms by the negative-coordinate set \(N=\{i:\alpha_i<0\}\).
Each group extends normally to \(\Omega_N\times D\): on a given compact
set choose inner contour radii for coordinates in \(N\) and outer radii
for all other coordinates. The inner-to-compact and compact-to-outer
radius ratios are strictly less than one. Cauchy's estimate bounds the
resulting sum by products of convergent geometric series, uniformly on
that compact set and on its compact set of \(z\)-parameters. This also
proves extension across newly allowed zero coordinates. The coefficient
projections and extensions commute with restriction in both the
projective charts and \(D\).

For each nonempty proper \(N\), choose \(j\notin N\). Extend every
\(N\)-component to the common domain \(\Omega_N\times D\), so that
restriction becomes the identity on its coefficient space. Regard
cochains as alternating functions of distinct vertices, zero on repeated
vertices and on faces not containing \(N\). Set \[
(hc)_{i_0\ldots i_{p-1}}=c_{j i_0\ldots i_{p-1}}.
\] After inserting \(j\), reorder with its sign and restrict to the
desired face. The term deleting the inserted vertex is the identity;
each other insertion-deletion pair has opposite signs. Thus \(dh+hd=1\).
In degree zero the would-be augmented value is zero, because the
singleton \(\{j\}\) does not contain \(N\); the same identity therefore
contracts the unaugmented complex. These operators act on normally
convergent series, by the preceding estimates, rather than on formal
expressions alone. For \(N=\varnothing\), the finitely many nonnegative
exponents of total weight \(d\) give the ordinary simplex complex, with
its coefficient module in degree zero. For \(N=\{0,\ldots,n\}\), the
finitely many strictly negative exponents of weight \(d\) occur only on
the full face, in degree \(n\). Empty exponent sets contribute zero.
When \(n=0\), these two possibilities both occur in degree zero and
together give the single basis element for the trivial line bundle on a
point. This is precisely the algebraic monomial computation with
coefficient ring \(\mathcal O(D)\), proving (7). \(\square\)

\textbf{Proposition 6.6.} Let \(f:Y\to X\) be a projective morphism of
complex schemes, and let \(F\) be coherent. Then the natural maps \[
(R^qf_*F)^{\mathrm{an}}
\longrightarrow R^qf^{\mathrm{an}}_*F^{\mathrm{an}}
\tag{9}
\] are isomorphisms for every \(q\ge0\). In particular the sheaves on
the right are coherent.

\textbf{Proof.} First take \(X=\mathbf A^m_{\mathbf C}\),
\(Y=\mathbf P^n_X\), and write \(r\) for the projection. By Lemma 6.5
and the local description of higher direct images as the sheafifications
of cohomology on inverse images of opens, comparison is an isomorphism
for every \(\mathcal O(d)\). Algebraically
\(R^qr_*\mathcal O(d)=V_d^q\otimes\mathcal O_X\), by the earlier
relative projective-space calculation. Analytically the same equality
follows on the polydisc basis from (7).

For any coherent \(H\) on \(\mathbf P^n_X\), both sides of (9) vanish
for \(q>n\). Algebraically use the finite standard affine cover.
Analytically use the same standard cover over a small polydisc; all its
intersections are the Stein manifolds just described, and
\(H^{\mathrm{an}}\) is coherent on them. Choose a finite twist
presentation \[
0\to K\to E\to H\to0,
\qquad E=\bigoplus_j\mathcal O(d_j),
\tag{10}
\] with \(K\) coherent, by the earlier Serre generation theorem over the
affine base.

Descend on \(q\), simultaneously for every coherent \(H\). Assume
comparison in degree \(q+1\) is an isomorphism for all coherent modules.
In the two long exact direct-image sequences from (10), comparison for
\(E\) is already an isomorphism. At any analytic base stalk, lift the
boundary of an analytic degree-\(q\) element uniquely through degree
\(q+1\) for \(K\). Its image in degree \(q+1\) for \(E\) is zero, so
algebraic exactness lifts it to degree \(q\) for \(H\). The difference
has zero boundary and can be corrected using degree \(q\) for \(E\).
This proves surjectivity for every \(H\). In particular it proves
surjectivity for \(K\). If an algebraic element for \(H\) has zero
analytic image, its boundary is zero by the degree-\((q+1)\) injectivity
for \(K\). Lift it from \(E\), correct that lift with the just-proved
surjectivity for \(K\), and use injectivity for \(E\) to conclude it is
zero. This proves injectivity. Analytification is exact on the base, so
these are two exact sequences of stalk modules. The induction proves (9)
in this ambient case.

For a general affine base \(X=\operatorname{Spec}A\), embed it as a
closed subscheme \(k:X\hookrightarrow\mathbf A^m\). A projective
embedding gives a closed immersion
\(i:Y\hookrightarrow\mathbf P^n_X\hookrightarrow\mathbf P^n_{\mathbf A^m}\).
Put \(H=i_*F\). Exact closed-immersion pushforward and its compatibility
with analytification identify the ambient comparison just proved with
\(k_*\) applied to (9). Closed-immersion pushforward detects
isomorphisms on its source, so (9) follows. Covering an arbitrary base
by affine opens proves the result: the natural comparison maps agree on
overlaps, and higher direct images commute with open restriction. This
proof never replaces a coherent module by a finite globally free
resolution. \(\square\)

\subsubsection{Cohomological GAGA for every proper
scheme}\label{cohomological-gaga-for-every-proper-scheme}

\textbf{Theorem 6.7.} If \(X\) is proper over \(\mathbf C\) and \(F\) is
coherent, then \[
H^q(X,F)\longrightarrow H^q(X^{\mathrm{an}},F^{\mathrm{an}})
\tag{11}
\] is an isomorphism for every \(q\ge0\).

\textbf{Proof.} Let \(\mathcal P(F)\) mean that all maps (11) are
isomorphisms. It holds for zero and is invariant under isomorphisms. The
natural comparison maps commute with boundaries because analytification
is exact and the maps are the morphism of cohomological delta functors
extending the map on sections. Two out of three in a short exact
sequence follows from its two long exact cohomology sequences.

Use
\href{proper-morphisms-and-coherent-direct-images.html\#3-the-property-form-of-dévissage}{Coherence
of higher direct images under proper morphisms, Theorem 3.1}, whose full
ideal-filtration argument permits one generic-rank-one witness for each
integral closed support. For every integral closed subscheme
\(i:Z\hookrightarrow X\), choose the earlier Chow construction's
projective modification \(p:Y\to Z\). Both
\(Y\to\operatorname{Spec}\mathbf C\) and \(p\) are projective. A very
ample line bundle \(L\) for an absolute projective embedding is also
\(p\)-ample: its embedding together with \(p\) is a closed embedding
into \(\mathbf P^N_Z\), by the closed-graph argument. Relative Serre
vanishing provides \(a\) with \[
R^jp_*L^a=0\qquad(j>0).
\] Set \(G=p_*L^a\), a coherent module by projective finiteness. Over
the dense isomorphism open it is an invertible module. Consequently
\(i_*G\) has support exactly \(Z\), generic stalk annihilated by the
maximal ideal of \(\mathcal O_{X,\eta_Z}\), and generic dimension one
over \(\kappa(\eta_Z)\).

Proposition 6.6 identifies \(G^{\mathrm{an}}\) with
\(p^{\mathrm{an}}_*L^{a,\mathrm{an}}\) and makes all higher analytic
direct images zero. The algebraic and analytic Leray sequences therefore
collapse to natural isomorphisms \[
H^q(Z,G)=H^q(Y,L^a),\qquad
H^q(Z^{\mathrm{an}},G^{\mathrm{an}})=H^q(Y^{\mathrm{an}},L^{a,\mathrm{an}}).
\] Projective cohomological GAGA on \(Y\) identifies the right sides.
Naturality of Leray and of direct-image comparison identifies this with
(11) for \(i_*G\), using exact closed-immersion pushforward. The earlier
dévissage theorem now proves \(\mathcal P(F)\) for every coherent \(F\),
including modules on a nonreduced \(X\). \(\square\)

\subsubsection{Morphisms and extension
classes}\label{morphisms-and-extension-classes}

\textbf{Lemma 6.8.} For coherent \(F,G\) on a complex scheme, the
sheaves \(\mathcal Ext^q_X(F,G)\) are coherent and \[
\mathcal Ext^q_X(F,G)^{\mathrm{an}}
\cong\mathcal Ext^q_{X^{\mathrm{an}}}(F^{\mathrm{an}},G^{\mathrm{an}})
\tag{12}
\] for every \(q\ge0\).

\textbf{Proof.} On an algebraic affine open, repeatedly present the
finitely generated syzygies of \(F\) by finite free modules.
Noetherianness ensures that this can be done to any prescribed finite
length. For a fixed \(q\), the finite free prefix through degrees
\(q+1\) computes \(\operatorname{Ext}^q\) by taking Hom into \(G\), then
a kernel modulo an image. The sheaf computation uses the same finite
prefix: free module sheaves are locally projective, so applying internal
Hom to it computes the corresponding sheaf Ext in these degrees. The
kernels and images are finite modules, hence coherent. Localization of
this finite computation is exact and commutes with its finite Hom terms,
so these local sheaves glue.

At a complex point, tensor the prefix with the flat analytic local
algebra \(B\) over the algebraic local ring \(A\). Exactness leaves a
free-resolution prefix for \(B\otimes_A F_x\). Finite freeness
identifies each analytic Hom term with its algebraic Hom term tensored
with \(B\); flatness preserves the kernels, images and quotient. This
proves (12) at every analytic stalk. The construction is natural and
agrees on overlaps, giving the sheaf isomorphism. \(\square\)

\textbf{Proposition 6.9.} For proper \(X\) and coherent \(F,G\), the
natural maps \[
\operatorname{Ext}^n_X(F,G)\longrightarrow
\operatorname{Ext}^n_{X^{\mathrm{an}}}(F^{\mathrm{an}},G^{\mathrm{an}})
\tag{13}
\] are isomorphisms for every \(n\ge0\). In degree zero this gives full
faithfulness; in degree one every analytic extension between algebraic
coherent modules is algebraic.

\textbf{Proof.}
\href{../derived-categories-and-sheaf-operations/internal-derived-hom-and-ext-sheaves.html\#3-adjunction-local-maps-and-ext}{Internal
derived Hom and Ext sheaves, Theorem 3.1 and equations (3.2)--(3.3)},
identifies global Ext with the hypercohomology of
\(K=R\mathcal Hom(F,G)\). Its cohomology sheaves are
\(\mathcal Ext^q(F,G)\), zero for \(q<0\). Applying the earlier
first-quadrant hypercohomology exact-couple construction to the good
truncations of \(K\) gives \[
E_2^{p,q}=H^p(X,\mathcal Ext^q(F,G))
\Longrightarrow\operatorname{Ext}^{p+q}_X(F,G).
\tag{14}
\] Here is why precisely the same construction applies: represent \(G\)
by an injective resolution in nonnegative degrees and take internal Hom
from \(F\). This is a nonnegative complex representing \(K\). Its
good-truncation triangles have successive terms
\(\mathcal Ext^q(F,G)[-q]\); their exact couples give (14), exactly as
in the earlier Leray proof. In total degree \(n\), only \(p,q\ge0\) with
\(p+q=n\) occur, so the abutment filtration has finitely many terms. No
finite projective dimension is asserted.

Exact analytification gives the natural derived comparison
\(K^{\mathrm{an}}\to R\mathcal Hom(F^{\mathrm{an}},G^{\mathrm{an}})\).
One construction is to analytify the injective-resolution Hom complex
and then map to the Hom complex of an analytic injective resolution; the
tensor-Hom evaluation map supplies the comparison, and injective
comparison makes it independent of choices. The finite local calculation
in Lemma 6.8 identifies the maps on cohomology sheaves. Naturality of
the truncation exact couples gives a morphism from (14) to its analytic
counterpart. Theorem 6.7 and Lemma 6.8 make every map on the \(E_2\)
page an isomorphism. Therefore every subsequent page, and every
successive quotient of the finite abutment filtration, is an
isomorphism. Induction up that finite filtration proves (13).

For clarity, \(\operatorname{Ext}^1\) here classifies short exact
sequences with specified ends. Embed \(G\) in an injective \(I\), with
quotient \(Q\). A map \(F\to Q\) yields an extension by pulling back
\(I\to Q\). Two such maps give equivalent extensions exactly when their
difference lifts to \(I\). Conversely a map from \(G\) to \(I\) extends
across any extension's middle module by injectivity and gives such a map
to \(Q\). Thus extensions are the cokernel of
\(\operatorname{Hom}(F,I)\to\operatorname{Hom}(F,Q)\), which is the
injective-resolution \(\operatorname{Ext}^1\). An extension of coherent
modules is coherent, algebraically and analytically. This construction
also shows that (13) sends an algebraic extension to its
analytification. \(\square\)

\subsubsection{Algebraizing an arbitrary coherent analytic
module}\label{algebraizing-an-arbitrary-coherent-analytic-module}

\textbf{Lemma 6.10 (a uniform support power).} Let \(X\) be proper,
\(Z\subset X\) a closed reduced subscheme with coherent ideal \(J\), and
\(M\) a coherent analytic module whose support is contained in
\(Z^{\mathrm{an}}\). Some power \((J^{\mathrm{an}})^r\) annihilates
\(M\).

\textbf{Proof.} At a point \(x\), the finite analytic module \(M_x\) has
support equal to the zero germ of its annihilator. The analytic
Nullstellensatz implies that every germ in \(J^{\mathrm{an}}_x\), whose
zero germ contains that support, lies in
\(\sqrt{\operatorname{Ann}(M_x)}\). Choose finitely many generators
\(g_i\) of that ideal. If \(g_i^{a_i}\) annihilates \(M_x\), every
product of \(\sum_i(a_i-1)+1\) generators contains one such power. Thus
a power of the whole ideal annihilates the stalk. Its product with \(M\)
is coherent, so its vanishing at this stalk holds on a neighbourhood.
Lemma 6.4 makes \(X^{\mathrm{an}}\) compact; choose finitely many of
these neighbourhoods and take their maximum exponent. This is the
required uniform \(r\). \(\square\)

\textbf{Theorem 6.11 (proper GAGA, with nilpotents).} For every proper
complex scheme \(X\), analytification is an equivalence \[
\operatorname{Coh}(X)\simeq\operatorname{Coh}(X^{\mathrm{an}}).
\] Together with Theorem 6.7 it preserves all coherent cohomology
groups.

\textbf{Proof.} Full faithfulness is Proposition 6.9 in degree zero.
Prove essential surjectivity by induction on \(\dim X\), proving the
reduced case first at each dimension, then all its nilpotent
thickenings. The empty case is immediate. Assume algebraization already
proved for all proper schemes of smaller dimension, including their
nilpotents, and first let \(X\) be reduced of the current dimension.

Choose Lemma 6.4's \(p:Y\to X\), which is an isomorphism over \(U\)
containing all generic points. For a coherent analytic module \(M\), its
analytic pullback \(p^{\mathrm{an},*}M\) is coherent, by finite
presentations. Since \(Y\) is projective, the already proved projective
GAGA algebraizes it to a coherent \(H\) on \(Y\), with a specified
isomorphism \(H^{\mathrm{an}}\cong p^{\mathrm{an},*}M\). Proposition 6.6
and projective finiteness give \[
N=p_*H\in\operatorname{Coh}(X),\qquad
N^{\mathrm{an}}\cong p^{\mathrm{an}}_*p^{\mathrm{an},*}M.
\] The pullback-pushforward adjunction has its canonical unit \[
u:M\longrightarrow N^{\mathrm{an}}.
\tag{15}
\] It is an isomorphism over \(U^{\mathrm{an}}\). Its coherent kernel
\(K\) and cokernel \(C\) are therefore supported on the analytification
of the proper algebraic closed subset \(Z=X\setminus U\). Give \(Z\) its
reduced structure, with ideal \(J\). Every component of \(Z\) has
smaller dimension than \(X\), because \(U\) contains every generic
point. Lemma 6.10 gives a common \(r\) with
\((J^{\mathrm{an}})^rK=(J^{\mathrm{an}})^rC=0\). Consequently \(K,C\)
are coherent modules on the analytic nilpotent thickening
\(Z_r=V(J^r)\), a proper scheme of smaller dimension. Induction
algebraizes them; exact closed-immersion pushforward supplies coherent
\(K_0,C_0\) on \(X\) with specified identifications
\(K_0^{\mathrm{an}}=K\), \(C_0^{\mathrm{an}}=C\).

Full faithfulness lifts the analytic quotient map
\(N^{\mathrm{an}}\to C\) to a map \(N\to C_0\). Its cokernel has zero
analytification and is therefore zero, so it is surjective. Let \(I_0\)
be its coherent kernel. Exactness identifies \(I_0^{\mathrm{an}}\) with
the image of (15). The given analytic short exact sequence \[
0\to K_0^{\mathrm{an}}\to M\to I_0^{\mathrm{an}}\to0
\] has a class in
\(\operatorname{Ext}^1_{X^{\mathrm{an}}}(I_0^{\mathrm{an}},K_0^{\mathrm{an}})\).
Proposition 6.9 lifts that class to an algebraic extension \[
0\to K_0\to F\to I_0\to0.
\] The middle module is coherent. Equality of extension classes gives an
analytic isomorphism \(F^{\mathrm{an}}\cong M\) compatible with the
specified ends. This proves the reduced case at the current dimension.

Finally let \(X\) be nonreduced of that dimension. Its coherent
nilradical \(A\) has \(A^s=0\) for some \(s\): on each of a finite
affine cover choose finitely many nilpotent generators, and take a
uniform sufficiently large power. The finite filtration \[
M\supset A^{\mathrm{an}}M\supset\cdots\supset(A^{\mathrm{an}})^sM=0
\] has coherent quotients annihilated by \(A^{\mathrm{an}}\). They are
coherent modules on \(X_{\mathrm{red}}^{\mathrm{an}}\), which algebraize
by the reduced case just proved at the same dimension. Regard their
algebraizations as modules on \(X\) by closed-immersion pushforward.
Starting with the bottom quotient, lift each successive extension using
Proposition 6.9 for this proper \(X\); the finite upward induction
algebraizes \(M\). This completes the dimension induction for all proper
schemes. \(\square\)

With identifications to the analytic module fixed, the algebraization is
unique up to the unique compatible isomorphism, by full faithfulness.
This assertion does not erase the automorphisms of a coherent module.

\subsection{7. Exercises with solutions}\label{exercises-with-solutions}

\textbf{Exercise 7.1 (easy).} Show why the cohomological comparison
theorem fails for \(\mathbf A^1\).

\textbf{Solution.} In degree zero it is the inclusion
\(\mathbf C[z]\to\mathcal O(\mathbf C)\). The entire function \(e^z\) is
absent from the image, as its derivatives never become zero. Thus even
degree-zero surjectivity fails. The projective hypothesis supplies the
compactness and twist arguments used above.

\textbf{Exercise 7.2 (medium).} Classify analytic line bundles on
projective space, explaining the dimension-zero exception.

\textbf{Solution.} Essential surjectivity algebraizes the bundle; local
faithful flatness descends its rank-one freeness as in Proposition 6.3.
A rational section and homogeneous prime-divisor equations make its
divisor a multiple of a hyperplane, so it is
\(\mathcal O(d)^{\mathrm{an}}\). For \(n\geq1\), the global-section
dimension distinguishes the trivial twist and hence makes \(d\) unique.
For \(n=0\), every \(\mathcal O(d)\) is the same one-dimensional bundle,
so there is no unique integer.

\textbf{Exercise 7.3 (medium).} Deduce Chow's theorem for a nonempty
analytic hypersurface of \(\mathbf P^n\) directly from its ideal, using
GAGA.

\textbf{Solution.} An analytic hypersurface in projective space has an
invertible coherent ideal: the convergent local rings are factorial, by
\href{../complex-analytic-spaces-and-coherent-sheaves/the-local-ring-of-holomorphic-germs.html\#3-unique-factorization}{The
local ring of holomorphic germs, Theorem 3.1}, so a reduced pure
codimension-one vanishing ideal is locally the product of its prime
equations. Algebraize this ideal and its inclusion into the structure
sheaf. The ideal descends to an algebraic line bundle by the same local
argument as in Proposition 6.3, so is \(\mathcal O(-d)\). Its inclusion
is a nonzero section of \(\mathcal O(d)\), hence a homogeneous
polynomial of degree \(d\). A nonempty zero set requires \(d>0\). The
ideal's analytic extension is the original vanishing ideal; radicality
descends by faithful flatness. Thus this polynomial defines exactly the
given reduced hypersurface; the ideal formulation handles reducible
hypersurfaces as well.

\textbf{Exercise 7.4 (medium).} Describe holomorphic maps
\(\mathbf P^1(\mathbf C)\to\mathbf P^1(\mathbf C)\).

\textbf{Solution.} Theorem 6.2 makes the map algebraic. Its pullback of
\(\mathcal O(1)\) is \(\mathcal O(d)\), and the pulled-back coordinate
sections generate it, so \(d\geq0\). They are homogeneous degree-\(d\)
polynomials \(F_0,F_1\) without a common projective zero. The map is
\([T_0:T_1]\mapsto[F_0(T):F_1(T)]\), and in an affine target chart it is
their rational quotient. For \(d=0\) it is constant; for \(d>0\) its
degree is \(d\).

\textbf{Exercise 7.5 (medium).} Compute analytic cohomology of
\(\mathcal O(d)\) on \(\mathbf P^1\), and verify comparison explicitly.

\textbf{Solution.} Use \(z=T_1/T_0\), with trivializations
\(e_0=T_0^d\), \(e_1=T_1^d=z^de_0\). The acyclic two-chart cover gives
the degree-one quotient of holomorphic functions on \(\mathbf C^*\) by
entire functions in \(z\) and \(z^d\) times entire functions in
\(z^{-1}\). In the Laurent expansion the first family removes exponents
\(a\geq0\); the second removes exponents \(a\leq d\). The remaining
basis is \(z^{d+1},\ldots,z^{-1}\) when \(d\leq-2\), and is empty
otherwise. Thus \(h^1=\max(-d-1,0)\). Global sections must have
exponents both \(a\geq0\) and \(a\leq d\), giving \(1,z,\ldots,z^d\) for
\(d\geq0\) and zero otherwise. Therefore \(h^0=\max(d+1,0)\); higher
cohomology vanishes. The algebraic Laurent Čech calculation has the
identical bases, and comparison takes each basis monomial to itself.

\textbf{Exercise 7.6 (harder).} State and prove the uniqueness of
algebraization with a fixed analytic identification. Give a
counterexample to uniqueness of the bare isomorphism.

\textbf{Solution.} For \(\phi:F^{\mathrm{an}}\to M\) and
\(\psi:G^{\mathrm{an}}\to M\), full faithfulness lifts \(\psi^{-1}\phi\)
uniquely to \(u:F\to G\). Lifting the inverse and using faithfulness
proves \(u\) is an isomorphism. It is the unique one compatible with the
identifications. On a nonempty projective variety, multiplication by any
\(\lambda\in\mathbf C^*\) is an automorphism of \(\mathcal O_X\); there
are many isomorphisms of that bare sheaf to itself. The compatibility
condition singles out the correct one.

\subsection{8. Sources and the comparison
chain}\label{sources-and-the-comparison-chain}

J.-P. Serre's
\href{https://www.numdam.org/item/AIF_1956__6__1_0/}{\emph{Géométrie
algébrique et géométrie analytique}}, Annales de l'Institut Fourier 6
(1956), 1--42, is the historical source for these three theorems and
their applications. The linked edition is free to read; no permission to
copy or translate it is asserted. Chow's theorem is credited to W.-L.
Chow. The arguments here are independently written; Section 1 links the
analytic programme proofs and Demailly's freely accessible exposition.
Component authorship and reuse terms remain recorded in the respective
lessons.

Kiran S. Kedlaya's
\href{https://ocw.mit.edu/courses/18-726-algebraic-geometry-spring-2009/resources/mit18_726s09_lec22_gaga/}{\emph{GAGA},
updated 30 April 2009}, MIT OpenCourseWare 18.726, treats comparison,
full faithfulness, algebraization and Chow in Sections 3--5 and 8. Its
notes are reusable under
\href{https://ocw.mit.edu/pages/privacy-and-terms-of-use/}{CC BY-NC-SA
4.0}. They offer useful parallel reading, while this course retains its
explicit Laurent contraction, two-stage descending comparison induction
and treatment of noninjective hyperplane multiplication.

For a broader framework, Jack Hall's
\href{https://arxiv.org/abs/1804.01976v3}{\emph{GAGA theorems},
arXiv:1804.01976v3, 17 May 2022}, Theorem A and Example 9.2, treats
analytic GAGA for proper schemes. Theorem A requires coherence, finite
total coherent cohomology, detection at closed points and specified
local flatness and residue-field conditions. Its non-Noetherian
extensions have further hypotheses. This public author draft is further
reading; the proofs above are supplied in the programme.

The corresponding verified open AI Integrated Stacks Project treatments
are
\href{https://kokunoyumeto.github.io/stacks-zh-hans-cn/en/gaga.html\#gaga-section-proof-gaga-cohomology}{cohomological
comparison},
\href{https://kokunoyumeto.github.io/stacks-zh-hans-cn/en/gaga.html\#gaga-section-proof-gaga-fully-faithful}{full
faithfulness},
\href{https://kokunoyumeto.github.io/stacks-zh-hans-cn/en/gaga.html\#gaga-lemma-projective-analytic-global-generation}{analytic
global generation},
\href{https://kokunoyumeto.github.io/stacks-zh-hans-cn/en/gaga.html\#gaga-section-completion-gaga-essential-surjectivity}{essential
surjectivity}, and
\href{https://kokunoyumeto.github.io/stacks-zh-hans-cn/en/gaga.html\#gaga-theorem-chow-closed-analytic-projective}{Chow's
theorem}. Their reference texts retain GFDL 1.2 and contain AI-written
additions not reviewed by the official Stacks maintainers. Section 1
states the analytic finiteness and acyclicity inputs; Section 2 gives
the Laurent structure-sheaf calculation assuming those inputs. The
referenced chapter also invokes analytic finiteness in its generation
argument.

Using the analytic programme results in Section 1 and the full quotient
construction in the preceding lesson, Sections 2--6 establish comparison
and algebraicity for projective complex schemes, including nonreduced
schemes. The chain is local analytic algebra and faithful flatness;
analytic Čech cohomology and finiteness; comparison for twists;
comparison for coherent sheaves; internal Hom and full faithfulness;
analytic generation and two presentations; algebraicity of full ideals
and graphs. Section 6 extends the result to every proper complex scheme:
relative projective comparison and rank-one dévissage give cohomological
comparison, local Ext and its finite abutment filtration give extension
comparison, and the modification unit reconstructs a coherent analytic
module from its two defects on a smaller closed support. A final finite
nilpotent filtration handles nonreduced schemes.
\end{document}
